\chapter{Introduction and Project Vision}
\label{chap:intro}

\section{Project Overview}
Video games represent one of the most prominent sectors of modern entertainment, increasingly relying on Procedural Content Generation (PCG) to maintain long-term replayability, mitigate content exhaustion, and reduce escalating production overhead. While early algorithmic generation depended on deterministic cellular automata, grammars, or Binary Space Partitioning (BSP), recent developments have advanced toward learned policies trained through Reinforcement Learning (PCGRL)~\cite{khalifa2020pcgrl}. 

However, existing PCGRL systems typically generate game content optimized for an abstract, aggregate statistical average or require human designers to author explicit target vectors manually~\cite{earle2021controllable}. When real human players interact with such systems, the static nature of generation inevitably conflicts with the broad distribution of human motor skill, tactical risk tolerance, spatial navigation pacing, and combat capabilities.

Project \textbf{FORGE} (\textit{Personalized Procedural Content Generation via Reinforcement Learning}, Project ID: F26-012-FORGE) is an online, closed-loop adaptive dungeon synthesizer engineered for 2D top-down combat roguelikes within the Godot 4.7 engine. Rather than generating content offline or relying on intrusive questionnaires, FORGE actively and passively records in-session gameplay telemetry (clearing time, death count, and damage-taken ratio). Together with the current difficulty settings and run depth, this telemetry forms an 8-dimensional observation that conditions a Proximal Policy Optimization (PPO)~\cite{schulman2017proximal} actor policy that outputs a continuous 8-dimensional difficulty action. The action sets enemy combat attributes (health, damage, speed, aggro radius), hazard density, resource scarcity, room count and an archetype-complexity value that decides the mix of five enemy types on the next floor. 

By pairing an offline Python training pipeline with a lightweight pure-GDScript neural network forward pass (well under 1\,ms, run once per floor), FORGE aims to sustain player engagement without perceptible difficulty thrashing; whether it does so for human players is evaluated in FYP-2.

\section{Problem Statement}
Procedural level generation in combat dungeon crawlers faces a critical trade-off between infinite variety and calibrated challenge:
\begin{enumerate}
    \item \textbf{The One-Size-Fits-All Failure Mode:} Content tuned for a nominal player inevitably under-challenges highly skilled veterans (inducing boredom and disengagement) and over-penalizes novice players (inducing anxiety, frustration, and premature churn).
    \item \textbf{The Brittle Rubber-Banding Fallacy:} Traditional Dynamic Difficulty Adjustment (DDA) systems adjust game parameters reactively using hardcoded threshold triggers (such as scalar health or damage multipliers). These systems are readily transparent to players, produce jarring swings in difficulty, and undermine player agency by penalizing proficient play.
    \item \textbf{The Limited Reach of Player-Adaptive PCGRL:} Most PCGRL generators are conditioned on random seeds or designer-set targets~\cite{earle2021controllable, earle2024pcgrlplus}. Player-adaptive RL generators do exist, notably experience-driven PCG via RL~\cite{shu2021edrl} and online player-adaptive generation~\cite{wang2022oparl}, and player-model-driven personalization predates RL~\cite{shaker2010personalized}. These systems target level geometry in platformers and are evaluated with simulated or artificial players; adapting several combat and environment attributes at once inside a shipped action game remains less explored.
\end{enumerate}

\subsection*{Central Research Question}
To address these interrelated deficiencies, Project FORGE is guided by the following central research question:
\begin{quote}
\textit{``Can a reinforcement-learning-based procedural content generation system utilize continuous in-session gameplay telemetry to infer an individual player's behavioral profile and synthesize solvable dungeon content whose multi-dimensional combat and environmental difficulty is measurably better matched to that player's observed capability than non-personalized procedural baselines?''}
\end{quote}

\section{Problem Context and Motivation}
A common psychological framing for game difficulty is Csikszentmihalyi's theory of flow~\cite{csikszentmihalyi1990flow}, whose application to digital games is reviewed in~\cite{michailidis2018flow}. Flow represents an optimal mental state characterized by deep absorption and intrinsic motivation, achieved when the challenges presented by an activity precisely balance the actor's perceived skills. When challenges outstrip capability, anxiety ensues; conversely, when skill dominates challenge, boredom terminates engagement.

In action roguelikes (e.g., \textit{Hades}, \textit{The Binding of Isaac}, \textit{Enter the Gungeon}), per-run replayability is paramount. However, players possess vastly different learning curves. A rigid difficulty curve leads to asymmetric retention: novices abandon the game within the first few floors, while experts speedrun through content without experiencing tactical resistance. 

Attempts by commercial studios to mitigate this via reactive DDA (such as the hidden difficulty rank reported for Capcom's \textit{Resident Evil 4}~\cite{engadget2015re4} or Valve's AI Director in \textit{Left 4 Dead}~\cite{booth2009l4d}) highlight the industry's demand for adaptive pacing. However, heuristic rule sets often create artificial difficulty adjustments that break immersion. There is an acute engineering and research demand for an AI system that treats difficulty as an organic, multi-attribute parameter space—incorporating adversary speed, aggro geometry, hazard density, and tactical archetype composition—learned and navigated via deep reinforcement learning.

\section{Project Vision}
The overarching vision of Project FORGE is to establish a unified, open-source, reproducible paradigm for player-conditioned procedural generation in action games. FORGE envisions a gameplay experience where:
\begin{itemize}
    \item Game content seamlessly conforms to the player's evolving flow channel without requiring manual difficulty selection menus (e.g., Easy, Normal, Hard).
    \item The adaptation mechanism operates unobtrusively, preserving the principle of ``change blindness'' so that players perceive the world as challenging yet fair, attributing their survival to personal mastery rather than patronizing AI assistance.
    \item The generative system runs deterministically on client hardware with no cloud or external server dependencies, and every generated floor is connected by construction (start-to-exit reachability is additionally checked by a regression test in CI).
\end{itemize}

\section{Project Objectives}
To realize this vision, the project team formulated six specific, measurable engineering and research objectives:
\begin{itemize}
    \item \textbf{OBJ-1 (Continuous Telemetry \& Profiling Engine):} Construct an unobtrusive telemetry logging framework capturing fine-grained in-session performance metrics (room traversal velocity, hit-to-damage ratio, trap trigger frequency, death count) at 60\,FPS without frame drops.
    \item \textbf{OBJ-2 (Reinforcement Learning Difficulty Policy):} Train a stable PPO policy in a simulated multi-room Gymnasium environment capable of mapping 8-dimensional observation states (three player-performance signals, four current difficulty settings and run depth) into smooth 8-dimensional difficulty parameter updates (seven numeric settings plus the archetype mix).
    \item \textbf{OBJ-3 (Two-Stage Decoupled Dungeon Synthesizer):} Implement a modular generator separating spatial topology generation from combat encounter placement, ensuring topological reachability (connected by construction and checked by a breadth-first search over the door graph on 500 CI seeds for each room count from 3 to 9).
    \item \textbf{OBJ-4 (Multi-Archetype Combat Dynamics):} Develop five distinct enemy archetypes (Melee Crawler, Ranged Pyromancer, Heavy Brute, Agile Skirmisher, Swarmer) with telegraphed attacks, and active player defensive mechanics (dodge roll).
    \item \textbf{OBJ-5 (Automated Continuous Integration Solvability Testing):} Establish a headless CI/CD pipeline using GitHub Actions that compiles Godot scripts, executes a headless reachability regression test over 3,500 seeded layouts, runs headless GDScript tests, and verifies policy export integrity on every push to the main branches and every pull request.
    \item \textbf{OBJ-6 (Empirical \& Psychometric Validation):} Establish a formal verification framework evaluating both objective performance metrics (solvability rate, difficulty error rate $\text{DER}$ = mean $|\bar{t}_{room} - 7.5\,\text{s}|/7.5\,\text{s}$, inference latency) and standardized psychometric instruments (Game Engagement Questionnaire, Player Experience of Need Satisfaction).
    \item \textbf{OBJ-7 (Learned and Explainable Difficulty Judge, added for Final-1):} Replace the hand-weighted difficulty formula with a Judge learned from telemetry and ratings (target MAE $\le 0.15$ on a $[0,1]$ scale), explain it with SHAP, and explain the generator's own choices with Policy SHAP.
\end{itemize}

\section{Project Scope}\label{sec:scope}
To ensure rigorous progress while avoiding scope creep, the project scope is formally partitioned across project milestones as shown in Table~\ref{tab:scope_matrix}.

\begin{table}[!htbp]
\centering
\caption{FORGE Project Scope Boundary and Delivery Matrix (updated for Final-1)}
\label{tab:scope_matrix}
\scriptsize
\begin{tabular}{p{0.23\textwidth}p{0.23\textwidth}p{0.23\textwidth}p{0.22\textwidth}}
\toprule
\textbf{Done (Mid-1 and Final-1)} & \textbf{Planned for Final-1 (not yet done)} & \textbf{Planned for FYP-2} & \textbf{Out of scope} \\
\midrule
Pixel-art sprites with animations, camera shake, damage numbers & Player model MLP1 trained on real telemetry & In-run equipment shop and currency (Mid-2) & 3D graphics, VR/AR \\
Five enemy archetypes and an 8-D action with archetype mix & Calibrated simulator (MLP1 + Judge in the reward) & Doors and keys in chests (Mid-2) & Online multiplayer \\
Dodge roll, five calibration floors, $CR_0$ & PPO retrained on the calibrated simulator & Judge retrained on 25+ players' ratings (Mid-2) & Procedural music or narrative \\
GDScript PPO forward pass, heuristic fallback & Policy SHAP on the PPO actor & Live SHAP panel inside Godot (Mid-2) & Training the policy on the player's machine \\
Static $+8\%$ baseline toggle & BSP baseline generator & 10-participant UAT with GEQ/PENS (Final-2) & Biometric or eye-tracking input \\
SQLite telemetry, consent screen, 1--5 floor rating & Layout-diversity metrics and one evaluation script & Integration and performance benchmarks (Mid-2) & \\
Offline LightGBM Judge with TreeSHAP (synthetic training) & Consent asked after Play (supervisor request, 18 Nov) & Public build & \\
Headless Godot CI (11 GDScript tests, SQLite test, 3,500 layouts) & & & \\
\bottomrule
\end{tabular}
\end{table}


\section{Project Modules}
The FORGE system is partitioned into five modules (Module 5 is new for Final-1), depicted conceptually in Figure~\ref{fig:modular_arch}:
\begin{enumerate}
    \item \textbf{Module 1: Telemetry \& Behavioral Profiling Engine (\texttt{TelemetryLogger.gd}):} Intercepts real-time gameplay events (player damage taken, death events, room entry/exit timestamps) and computes the three normalized player-performance signals (clearing time, deaths, damage-taken ratio) of the observation vector $s_t \in \mathbb{R}^8$; \texttt{DifficultyManager.gd} completes it with the four current difficulty settings and run depth. Computes the Initial Combat Rating ($CR_0$) during cold-start calibration and writes one row per floor to an SQLite database (\texttt{user://forge\_telemetry.db}) through \texttt{SqliteBackend.gd}, with a JSON fallback.
    \item \textbf{Module 2: RL Generator \& In-Engine Inference Engine (\texttt{PPOPolicy.gd}):} Ingests observation vectors and executes an embedded matrix forward pass of the trained actor network to output normalized action vectors ($a_t \in [-1, 1]^8$). Manages graceful degradation via rule-based heuristic fallback if weights are inaccessible.
    \item \textbf{Module 3: Dungeon Layout \& Encounter Synthesizer (\texttt{DungeonGenerator.gd}, \texttt{DifficultyManager.gd}):} Operates the two-stage generation pipeline. Synthesizes spatial room layouts, opens matching doors between adjacent rooms so that the layout is a connected graph, denormalizes PPO action parameters into concrete game bounds ($\vec{E}, \vec{M}$), and populates rooms with adversary groups and hazards.
    \item \textbf{Module 4: Core Roguelike Gameplay \& Presentation Engine (\texttt{Main.gd}, \texttt{Player.gd}, \texttt{Enemy.gd}, \texttt{DefenseOverlay.gd}):} Drives 60\,FPS real-time physics, combat hitboxes, projectile kinematics, health tracking, player dash mechanics, camera shake decay, floating combat indicators, and the supervisor defense inspection overlay, the consent screen and the post-floor difficulty rating.
    \item \textbf{Module 5: Judge, Player Model and Explanations (Python, offline; new for Final-1):} The learned Judge (\texttt{ml\_spike/judge/}), TreeSHAP Design Insight Reports, the planned player model MLP1, the calibrated simulator and Policy SHAP (Section~\ref{sec:revised_arch}).
\end{enumerate}

\begin{figure}[!htbp]
\centering
\begin{tikzpicture}[node distance=1.5cm, auto,
    block/.style={rectangle, draw=blue!70, fill=blue!10, rounded corners, text width=3.2cm, text centered, minimum height=1.2cm, font=\small\bfseries},
    data/.style={rectangle, draw=orange!70, fill=orange!10, text width=2.8cm, text centered, minimum height=1.0cm, font=\footnotesize},
    line/.style={draw, -latex, thick}]
    
    \node [block] (game) {Module 4: Core Gameplay Engine (Godot 4.7)};
    \node [block, right=3.0cm of game] (telem) {Module 1: Telemetry Profiling Engine};
    \node [block, below=1.3cm of telem] (rl) {Module 2: PPO Policy Inference};
    \node [block, below=1.3cm of game] (dungeon) {Module 3: Dungeon \& Encounter Generator};
    \node [data, above=0.9cm of game] (player) {Player (human)};
    \node [data, below=0.9cm of rl] (offline) {Offline PPO training (Python)};

    \path [line, <->] (player) -- node [right, font=\footnotesize] {controls $\leftrightarrow$ game view} (game);
    \path [line] (game) -- node [above, font=\footnotesize] {Combat events} (telem);
    \path [line] (telem) -- node [right, font=\footnotesize] {Observation $s_t$} (rl);
    \path [line] (rl) -- node [below, font=\footnotesize] {Action $a_t$} (dungeon);
    \path [line] (dungeon) -- node [left, font=\footnotesize] {Generated level} (game);
    \path [line, dashed] (offline) -- node [right, font=\footnotesize] {trained weights (JSON)} (rl);
\end{tikzpicture}
\caption{High-Level Modular Interaction Flow of Project FORGE}
\label{fig:modular_arch}
\end{figure}

\section{Stakeholders and User Classes}
The primary stakeholders and user classes interacting with FORGE are summarized in Table~\ref{tab:stakeholders}.

\begin{table}[!htbp]
\centering
\caption{Stakeholders and User Classification}
\label{tab:stakeholders}
\small
\begin{tabular}{p{0.22\textwidth}p{0.35\textwidth}p{0.35\textwidth}}
\toprule
\textbf{User Class / Stakeholder} & \textbf{Characteristics \& Needs} & \textbf{System Interaction Points} \\
\midrule
\textbf{End-User Player} & Casual or enthusiast action roguelike gamers seeking engaging, personalized challenge without tedious setup. & Gameplay controls (WASD, Mouse Sword Attack, Spacebar Dash), HUD, room traversal. \\
\textbf{Game Balance Designer} & Evaluators requiring transparency into how AI difficulty parameters adjust relative to telemetry. & Real-time Defense Overlay, parameter manual sliders, telemetry export logs. \\
\textbf{Academic Researchers} & Investigators evaluating PCGRL adaptability, solvability, and psychometric flow alignment. & Offline training scripts, policy weights JSON, automated headless CI test reports. \\
\textbf{Supervisory Panel} & Evaluators assessing project progress, code ownership, engineering rigor, and rubric compliance. & Form 2 Live Demonstration, Form 3 Written Mid-Report, Git commit evidence. \\
\bottomrule
\end{tabular}
\end{table}

\section{Success Criteria}
To provide objective benchmarks for project evaluation, the success criteria listed in Table~\ref{tab:success_criteria} were baselined.

\begin{table}[!htbp]
\centering
\caption{Measurable Project Success Criteria}
\label{tab:success_criteria}
\scriptsize
\begin{tabular}{p{0.07\textwidth}p{0.27\textwidth}p{0.24\textwidth}p{0.34\textwidth}}
\toprule
\textbf{ID} & \textbf{Success Metric} & \textbf{Target} & \textbf{Final-1 Status} \\
\midrule
SC-01 & Solvability (start-to-exit reachability) & $100\%$ of seeded layouts & Met: 3,500/3,500 layouts in CI (room counts 3--9) \\
SC-02 & In-engine inference latency & median $<1.0$\,ms per call & Met at Mid-1 (older policy); \tofill{re-measure shipped 8-output policy with \texttt{benchmark\_inference.gd}} \\
SC-03 & Frame rate & $\ge 60$\,FPS at $1280\times720$ & Now logged per floor (\texttt{frame\_p95\_ms}); \tofill{report p95 frame time from the 29 real floors} \\
SC-04 & Cold-start calibration & $CR_0 \in [10, 95]$ after 5 floors & Met (headless test \texttt{test\_cold\_start}) \\
SC-05 & Fault tolerance & No crash with missing weights & Met (heuristic fallback) \\
SC-06 & CI pass rate & All CI jobs green on \texttt{main} & \tofill{link or screenshot of the latest green CI run} \\
SC-07 & Difficulty error rate DER (pacing error) & $\le 0.15$ & Met in simulation for all six profiles ($0.12$--$0.13$); not yet measured on humans \\
SC-08 & Judge prediction error & MAE $\le 0.15$ on $[0,1]$, better than the linear formula & Met on synthetic players (MAE $0.052$); not met on 29 real floors (pilot) \\
SC-09 & Convergence time & Pacing error $\le 0.15$ within 3--4 floors after a skill change & Met in simulation: recovers in 2 floors after a $-0.30$ skill drop \\
SC-10 & Adaptation stability & Mean $|\Delta a|$ per floor $\le 0.35$ & Met in simulation: $0.058$ (worst profile $0.065$) \\
SC-11 & Layout diversity & FORGE at least as diverse as BSP & \stplan{} (needs BSP and diversity script) \\
SC-12 & Player experience & GEQ/PENS gain over baselines & Planned for FYP-2 UAT \\
\bottomrule
\end{tabular}
\end{table}


\section{Constraints and Assumptions}
\subsection*{Technical Constraints}
\begin{itemize}
    \item \textbf{Engine Constraints:} Development is confined to the Godot Engine 4.x; development, benchmarking and CI all use Godot 4.7.2. In-game forward passes must be coded in native GDScript to prevent external library or DLL installation failures on end-user machines.
    \item \textbf{Execution Budget:} Total frame compute time must not exceed $16.6$\,ms to maintain 60\,FPS vertical sync on standard hardware (Intel Core i5 / GTX 1060 equivalent).
    \item \textbf{Offline Training Horizon:} Policy training is constrained to CPU/GPU environments running Python 3.10 or newer (CI uses 3.12) with PyTorch and Stable-Baselines3.
\end{itemize}

\subsection*{Operational Assumptions}
\begin{itemize}
    \item Players interact using standard keyboard and mouse peripherals.
    \item Single-player execution without concurrent cloud dependency.
    \item Player performance changes gradually across floors, with sudden skill deviations limited in effect by the hard bounds on every difficulty parameter.
\end{itemize}

\section{Ethical, Legal and DEI Considerations}\label{sec:ethics_intro}
\subsection*{Data Privacy and Player Governance}
Adaptive difficulty systems profile players, so the data they collect must be handled with care. This section was rewritten for Final-1 because the Mid-1 text (``zero PII, nothing leaves the machine'') is no longer accurate:
\begin{itemize}
    \item \textbf{What is stored:} only gameplay numbers per floor (time, deaths, damage, settings, frame times, the optional 1--5 rating). No names, emails, IP addresses or device serial numbers are collected.
    \item \textbf{Pseudonymous ID:} the game stores a random \texttt{install\_id} that stays the same across sessions on one machine. It does not identify a person by itself, but it does link sessions, so the data is \emph{pseudonymous}, not anonymous.
    \item \textbf{Consent:} a consent screen asks whether the session may be used for research. Every floor is still written to the local database, but with \texttt{consent = 0}, and the Judge and all analysis use consented sessions only. The supervisor asked (18 Nov 2026) that consent be asked after the Play button; this change is planned.
    \item \textbf{Data leaving the machine:} the game has no network code. Playtesters send their \texttt{forge\_telemetry.db} file to the team by hand (\texttt{PLAYTEST.md}). \tofill{who holds the files, where they are stored, how long they are kept, and how a tester can ask for deletion}.
    \item \textbf{Ethics approval:} \tofill{confirm with the supervisor whether university ethics approval is needed before the 25+ player collection and the UAT}.
\end{itemize}

\subsection*{Game Addiction and Dark Pattern Safeguards}
Commercial engagement-driven DDA often employs predatory algorithms designed to artificially induce near-miss scenarios to maximize monetization or daily active time. FORGE explicitly rejects dark patterns: difficulty adaptation is oriented purely toward sustaining Csikszentmihalyi's flow channel for intrinsic player satisfaction, without manipulative churn mechanics or artificial loot throttling.

\subsection*{Diversity, Equity, and Inclusion (DEI)}
By calibrating combat difficulty to individual performance, FORGE aims to improve gameplay accessibility: novice players should meet fewer repeated, punitive failure walls, while expert players encounter escalating challenges. Whether this holds for players with motor impairments has not been tested and would require a dedicated study. Visual indicators (such as glowing yellow telegraph rings on Pyromancer fireballs) ensure tactical clarity for players with varying reaction times.
